\section{Conclusiones} \label{sec:conclusiones}
El objetivo de este trabajo era introducir, en primer lugar, el concepto de shell galáctica y exponer el principal modelo propuesto para su formación; además de alguna de sus propiedades como, en particular, la metalicidad. El siguiente paso consistía en comprobar si las observaciones confirmaban estas predicciones. Como caso de estudio, seleccionamos una galaxia situada a una distancia mucho mayor que la de otros trabajos similares y situada en un cúmulo, una circunstancia muy inusual. No obstante, aprovechando la resolución y las capacidades de observación en el infrarrojo del Telescopio James Webb, fuimos capaces de identificar la presencia de shells en el halo de la galaxia. El hecho de que no fuese posible realizar el mismo análisis morfológico con las imágenes del Telescopio Hubble subraya la importancia de los instrumentos más modernos en la identificación de objetos astrofísicos de bajo brillo superficial. Los datos de HST, junto con los del espectrógrafo MUSE fueron, no obstante, cruciales para realizar un análisis espectrofotométrico completo en múltiples bandas del rango óptico e infrarrojo cercano. Esto permitió realizar un ajuste completo de las medidas a un modelo para la distribución espectral de energía de una shell, que se comparó con el de una región de referencia sin presencia de shells. Las propiedades físicas obtenidas a partir de este modelo no se ajustaron a lo que esperábamos según las predicciones teóricas, pero el análisis señala un camino para futuras investigaciones en las que se aprovechen todas las herramientas que tenemos disponibles para la caracterización de este tipo de objetos.